% Source: https://artofproblemsolving.com/community/c3h1440649_easy_one
% Stable identity: aops-topic-1440649-q01
Fix a prime $p$, and let $\alpha$ and $\beta$ be the exponents of $p$ in $a$ and $b$. The divisibilities give, for every $k\ge1$,
\[
(2k-1)\alpha\le2k\beta\le(2k+1)\alpha .
\]
If $\alpha=0$, the right inequality with $k=1$ gives $2\beta\le0$, so $\beta=0=\alpha$. If $\alpha>0$, dividing by $2k\alpha$ gives
\[
1-\dfrac{1}{2k}\le\dfrac{\beta}{\alpha}\le1+\dfrac{1}{2k}\quad\text{for all }k,
\]
so $\left|\dfrac{\beta}{\alpha}-1\right|\le\dfrac{1}{2k}$ for every $k$, which forces $\dfrac{\beta}{\alpha}=1$. In either case $\alpha=\beta$. Since $a$ and $b$ have the same exponent of every prime, $a=b$. $\blacksquare$
